\documentclass[../main.tex]{subfiles}
\begin{document}
%%%%%%%%%%%%%%%%%%%%%%%%%%%%%%%%%%%%%%%%%%%%%%%%%%%%%%%%%%%%%%%%%%%%%%%%%%%%%%%%%%%%%%%%%%%%%%%%%%%
\chapter{Drawing Basics}

Everything you'll do in OmniGraffle involves drawing. Whether you're creating the basic shapes
for an organizational chart or prototyping a user interface for your next iOS or Mac app,
OmniGraffle is packed with tools to feed your creativity.

Before you dive in, though, take some time to learn some of the basics of drawing in
OmniGraffle. There are many tips and tricks here that can help you master OmniGraffle in no
time.

\section{Creating Shapes}

As mentioned in Basic Concepts, drawing in OmniGraffle takes the form of \textit{objects} on
layers of a canvas. Shapes are the most common type of OmniGraffle object you can add to a
project, and the most basic way to create, or \textit{draw}, a shape on the canvas is with the
\textbf{Shape tool}.

Select the Shape tool in the \textbf{Tool Palette}.

\ogfigure{0.800}{res/og7_drawingbasics01.png}

Next, open the \textbf{Object} inspectors by choosing \textbf{Inspectors} $\blacktriangleright$
\textbf{Object}. You can also open the Object inspectors by pressing \textbf{Command-1}. Each
of the inspector tabs can be accessed with a keyboard shortcut.

Locate the \textbf{Shape} inspector and \textbf{Option-click} on either its disclosure triangle
or on the inspector name. This opens the Shape inspector and collapses the other Object
inspectors in that tab.

\ogfigure{0.800}{res/og7_drawingbasics02.png}

Click the \textbf{Shape Selection} button; this opens a popover from which you can choose a
shape. \textbf{Double-click} on the square/rectangle shape in the upper-left corner of the
popover.

\ogfigure{0.800}{res/og7_drawingbasics03.png}

With the Shape tool and the Rectangle shape selected, drag across the canvas until you have a
rectangle that is the size you want. When you draw a shape, the mouse pointer changes to a set
of crosshairs.

\ogfigure{0.800}{res/og7_drawingbasics04.png}

\begin{ogtip}
  If you want to draw a perfect square, hold down the \textbf{Shift} key while drawing to restrain
  the dimensions of the shape.
\end{ogtip}

It doesn't have to be perfect right now; to set an exact size, use the \textbf{Geometry} object
inspector (\textbf{Command--1}).

\ogfigure{0.800}{res/og7_drawingbasics05.png}

The top two fields in the Geometry inspector let you know the position of the object in relation
to the canvas origin, while the middle two fields reflect the object's width and height. The
controls on the bottom row allow you to rotate the object or flip it horizontally or
vertically.

\section{Selecting, Resizing, and Moving Objects}

To select an object, choose the \textbf{Selection} tool
(\ogiconw{19.4pt}{res/og7_tools_selection.png}) from the Tool Palette, and then click the
object. When the object is highlighted, you will see eight handles along its edges, plus a
round handle at its center if you choose \textbf{View} $\blacktriangleright$ \textbf{Movement
Handles}.

\ogfigure{0.723}{res/og7_drawingbasics06.png}

If you have selected two or more objects, the area of the selected objects receives the eight
resizing handles. This allows you to position objects on the canvas first, and then later
resize select objects as needed.

\ogfigure{0.800}{res/og7_drawingbasics07.png}

To perfectly align two objects with each other, \textbf{click-and-drag} one of the objects into
position. As you drag the object into position, blue \textbf{Smart Guides} appear to help you
align the objects.

\ogfigure{0.800}{res/og7_drawingbasics31.png}

For more precise control over aligning objects, select the objects that you want to be aligned,
and use the \textbf{Alignment} inspector to align and specify the distance between objects
vertically or horizontally.

\ogfigure{0.800}{res/og7_drawingbasics32.png}

\begin{ognote}
  See Positioning Objects Neatly with the Alignment Inspector to learn more about the Alignment
  inspector.
\end{ognote}

If you have an object on the canvas and you'd like to change its shape, select the shape and
then use the \textbf{Shape} inspector's Shape Popover to select a different shape:

\ogfigure{0.800}{res/og7_drawingbasics33.png}

If you have two objects of different sizes on the canvas and you'd like them to be the same
size, do the following:

\begin{enumerate}
  \item Select the object that has the dimensions you'd like the others to emulate.
  \item \textbf{Command-click} on the objects that you want to be the same size as your first
    selection.
  \item Choose \textbf{Arrange} $\blacktriangleright$ \textbf{Size} and then choose one of the
    resizing options:

    \begin{itemize}
      \item Make Same Width
      \item Make Same Height
      \item Make Same Size
      \item Make Natural Size
      \item Size to Fit Image (this option is only available if you are trying to fit an image
        within the bounds of a shape)
    \end{itemize}
\end{enumerate}

\ogfigure{0.800}{res/og7_drawingbasics34.png}

Once you have selected an object, you can modify it with the \textbf{Inspectors} or use the menu
commands to alter its appearance.

\begin{itemize}
  \item To resize an object, first select the object and then drag one of its resizing handles
    until the object is the size you'd like.

    \ogfigure{0.800}{res/og7_drawingbasics08.png}
  \item To select or deselect more objects without losing the selection you've already made,
    hold \textbf{Command} as you click other objects.

    \ogfigure{0.800}{res/og7_drawingbasics09.png}
  \item To select several adjacent objects at once, click in an empty area of the canvas and
    drag a box around the objects that you want to select. If you hold down the \textbf{Option}
    key while dragging, you'll get only the objects that fit completely inside the selection
    rectangle.

    \ogfigure{0.800}{res/og7_drawingbasics10.png}
  \item To rotate an object, hold down the \textbf{Command} key while mousing over one of the
    eight handles. The pointer changes from a standard resizing control to a curved one. An
    angle control appears inside of the object, along with a helper dialog that shows the degree
    of the angle.

    \ogfigure{0.800}{res/og7_drawingbasics11.png}

    \begin{ogtip}
      If you are using a laptop or touch-sensitive trackpad, you can enable Multi-Touch trackpad
      gestures for use with OmniGraffle (see Multi-Touch). With Multi-Touch gestures enabled, you
      can use two fingers to rotate objects.
    \end{ogtip}
  \item To \textbf{duplicate} objects you've drawn, hold down the \textbf{Option} key and then
    click on and drag the objects. An exact duplicate appears on the canvas. Unlike using Copy
    and Paste from the \textbf{Edit} menu (which places the pasteboard content on the currently
    active layer), the duplicated objects remain on their original layers.
\end{itemize}

\ogfigure{0.769}{res/og7_drawingbasics12.png}

\begin{ognote}
  Duplicating objects in this way does not replace other copied content currently on the
  pasteboard.
\end{ognote}

\begin{itemize}
  \item To move an object, use the \textbf{Selection} tool to click and hold on the object, and
    then drag it to its new location. \textbf{Smart Alignment Guides} (the light blue lines)
    appear to help you align objects to one another, while \textbf{Smart Distance Guides} appear
    beneath the objects so you can space them evenly.
\end{itemize}

\ogfigure{0.800}{res/og7_drawingbasics13.png}

\subsection{Grouping and Ungrouping Objects}

A group behaves like a single object; you can select it, move it around, resize it, use
inspectors to change its style and properties, connect lines to it, and so on. There is no
limit to the number of objects that can be part of a group; you can even group multiple groups
together to create subgroups.

\begin{itemize}
  \item To group two or more objects together so they can be moved or styled as a single unit,
    select the objects and choose \textbf{Arrange} $\blacktriangleright$ \textbf{Group}
    (\textbf{Shift-Command-G}).

    \ogfigure{0.800}{res/og7_drawingbasics14.png}
  \item To select an object that's part of a group:

    \begin{itemize}
      \item Click once to select the group as a whole, then click again to select an individual
        member of the group. You can keep clicking to select members of groups within groups.
    \end{itemize}

    \ogfigure{0.800}{res/og7_drawingbasics14a.png}

    \begin{itemize}
      \item In the Sidebar, click the disclosure triangle next to the group to reveal the objects
        or subgroups it contains, and then select the object you need:
    \end{itemize}

    \ogfigure{0.800}{res/og7_drawingbasics14b.png}

    \begin{ogtip}
      Similarly, if you have OmniGraffle Pro, you can select the individual objects within a
      \textbf{Subgraph} or \textbf{Table} by opening their object group in the Sidebar. With a
      Subgraph or Table open in the Sidebar, you can drag other objects on the canvas into or out
      of a Subgraph or Table.
    \end{ogtip}
  \item To group two groups together (thus creating a subgroup), select the groups that you
    would like to group together\ldots{}

    \ogfigure{0.800}{res/og7_drawingbasics14c.png}

    \ldots{}and then choose \textbf{Arrange} $\blacktriangleright$ \textbf{Group}
    (\textbf{Shift-Command-G}):

    \ogfigure{0.800}{res/og7_drawingbasics14d.png}
  \item To add an object to an existing group, \textbf{select and drag} the object in the
    Sidebar\ldots{}

    \ogfigure{0.800}{res/og7_drawingbasics14e.png}

    \ldots{}and then drop it into the group:

    \ogfigure{0.800}{res/og7_drawingbasics14f.png}
  \item To take a group apart, select it and choose \textbf{Arrange} $\blacktriangleright$
    \textbf{Ungroup} (\textbf{Shift-Command-U}).
\end{itemize}

\section{Styling Shapes with the Object Inspectors}

As you will learn more in the Diagramming Basics tutorial, you can apply styles to the shapes
you create with the \textbf{Object inspectors}, available on the right side of the canvas.

To style an object:

\begin{enumerate}
  \item Select the object you want to style with the \textbf{Selection} tool:

    \ogfigure{0.800}{res/og7_drawingbasics24.png}

    When you select the object, the Inspectors sidebar conveniently switches to the Object
    inspectors.
  \item \textbf{Option-click} on the \textbf{Fill} inspector. This opens the Fill inspector, and
    collapses the other Object inspectors in the sidebar:

    \ogfigure{0.800}{res/og7_drawingbasics25.png}
  \item Click one of the \textbf{Color Swatches} or use the \textbf{Fill Color} menu to apply a
    fill color to the object you selected on the canvas:

    \ogfigure{0.800}{res/og7_drawingbasics26.png}

    \ogfigure{0.800}{res/og7_drawingbasics27.png}
  \item Similarly, you can use the \textbf{Stroke} inspector to change the point size of the
    line around the object and its color:

    \ogfigure{0.800}{res/og7_drawingbasics28.png}

    \ogfigure{0.800}{res/og7_drawingbasics29.png}
\end{enumerate}

\subsection{Connecting Shapes with Lines}

You can connect any two objects with the \textbf{Line} tool
(\ogiconw{19.4pt}{res/og7_tools_line.png}).

\begin{enumerate}
  \item Select the Line tool.

    \ogfigure{0.800}{res/og7_drawingbasics15.png}
  \item Click the first object you would like to connect; this object becomes the
    \textit{source} of the line. You'll notice that when you move the mouse over the object,
    that a blue pulsing border appears around the object.

    \ogfigure{0.800}{res/og7_drawingbasics16.png}
  \item Click the second object you would like to connect; this object becomes the
    \textit{destination} of the line. When you drag the line from the first object to the next,
    you'll notice that the blue pulsing border appears around the object you're dragging the
    line to.

    \ogfigure{0.800}{res/og7_drawingbasics17.png}
  \item When the objects are connected with a line, they stay connected no matter how much you
    change or move them around. That means the logical integrity of your diagram is preserved
    regardless of how you choose to style the objects.

    \ogfigure{0.800}{res/og7_drawingbasics18.png}
  \item You can disconnect a line from an object by selecting the line and then dragging one of
    its end points. To delete the line, just select the line and press \textbf{Delete}.

    \ogfigure{0.800}{res/og7_drawingbasics19.png}
  \item You can add points to a line to control its routing by double-clicking the line at the
    place where you would like to add the point, and dragging the point to reposition the line.
    This behavior can be changed in Drawing Tools preferences.

    \ogfigure{0.800}{res/og7_drawingbasics35.png}
\end{enumerate}

\begin{ogtip}
  To create an orthogonal line like the one in the figure above, click the line to select it and
  use the Line inspector to change the line type to \textbf{Orthogonal}.
\end{ogtip}

\subsection{Styling Connecting Lines and Adding Line Ends}

The connecting lines that you draw between objects on the canvas can be styled, just like any
other object you create in OmniGraffle. For example, you can change the point size and color
of the line with the \textbf{Stroke} inspector:

\ogfigure{0.800}{res/og7_drawingbasics20.png}

And, if you are creating a flowchart or diagram where showing the flow of a process is needed,
you can add line ends to the \textit{tail} or \textit{head} of the line. Using the
\textbf{Line} inspector, you can add a \textbf{Line Tail} or \textbf{Line Head}, such as
circles, squares, and various sets of arrows, depending on what it is you are diagramming.

\ogfigure{0.800}{res/og7_drawingbasics21.png}

The various line ends in OmniGraffle have some smart behavior. When you change the point size of
the line---say, from 1 pt to 2 pts---the arrow scales with the point size of the line.

\ogfigure{0.800}{res/og7_drawingbasics22.png}

However, if all you want is a bigger arrow, you can adjust the size of the arrow using the scale
field beneath the line tail or head pop-up menus. The scale fields accept positive and
negative numbers, so you can have teeny-tiny or really huge arrows.

\ogfigure{0.800}{res/og7_drawingbasics23.png}
%%%%%%%%%%%%%%%%%%%%%%%%%%%%%%%%%%%%%%%%%%%%%%%%%%%%%%%%%%%%%%%%%%%%%%%%%%%%%%%%%%%%%%%%%%%%%%%%%%%
\end{document}
